% =============================================================================
% ForensiQ — Levantamento de Requisitos (MoSCoW)
% UC 21184 · Projeto de Engenharia Informática · UAb
% =============================================================================
\documentclass[a4paper,11pt]{article}

\usepackage[utf8]{inputenc}
\usepackage[T1]{fontenc}
\usepackage[portuguese]{babel}
\usepackage[top=2.5cm, bottom=2.5cm, left=3cm, right=2.5cm]{geometry}
\usepackage{booktabs}
\usepackage{longtable}
\usepackage{array}
\usepackage{parskip}
\usepackage{hyperref}
\usepackage{enumitem}

\hypersetup{
    colorlinks=true,
    linkcolor=black,
    urlcolor=blue,
    pdfauthor={João Rodrigues},
    pdftitle={ForensiQ — Levantamento de Requisitos}
}

\newcolumntype{L}[1]{>{\raggedright\arraybackslash}p{#1}}

% =============================================================================
\begin{document}
% =============================================================================

\begin{center}
    {\LARGE\bfseries Levantamento de Requisitos}\\[0.8em]
    {\large ForensiQ --- Plataforma Modular de Gestão de Prova Digital\\
    para First Responders}\\[1.2em]
    \begin{tabular}{ll}
        \textbf{Versão:}    & 1.0 · 22 março 2026 \\
        \textbf{Referência MoSCoW:} &
            \href{https://www.productplan.com/glossary/moscow-prioritization/}%
                 {productplan.com/glossary/moscow-prioritization} \\
    \end{tabular}
\end{center}

\vspace{0.5em}
\hrule
\vspace{1.5em}

% =============================================================================
\section{Método MoSCoW}
% =============================================================================

\begin{longtable}{L{3cm} L{10cm}}
    \toprule
    \textbf{Categoria} & \textbf{Significado} \\
    \midrule
    \endhead
    \textbf{Must have}   & Obrigatório. Sem isto o projecto não é entregável. \\[0.3em]
    \textbf{Should have} & Importante mas não crítico. Incluir se o tempo permitir. \\[0.3em]
    \textbf{Could have}  & Desejável. Só se tudo o resto estiver concluído. \\[0.3em]
    \textbf{Won't have}  & Explicitamente fora do âmbito desta versão. \\
    \bottomrule
\end{longtable}

% =============================================================================
\section{Requisitos funcionais}
% =============================================================================

\subsection*{Must have}

\begin{itemize}[nosep, leftmargin=*]
    \item RF01 --- Autenticação JWT com perfis agente e perito forense digital
    \item RF02 --- Criação e gestão de ocorrências com georreferenciação automática
    \item RF03 --- Registo de prova com fotografia, metadados GPS, timestamp e hash SHA-256
    \item RF04 --- Módulo de forense digital: ficha completa de dispositivo digital
                   (tipo, marca, modelo, estado, número de série)
    \item RF05 --- Cadeia de custódia com máquina de estados e log \emph{append-only} imutável
    \item RF06 --- Exportação de relatório de ocorrência em PDF
    \item RF07 --- API REST com mínimo 10 endpoints documentados via Swagger UI
\end{itemize}

\subsection*{Should have}

\begin{itemize}[nosep, leftmargin=*]
    \item RF08 --- Dashboard de estado das provas por ocorrência
    \item RF09 --- Pesquisa e filtragem de ocorrências e provas
\end{itemize}

\subsection*{Could have}

\begin{itemize}[nosep, leftmargin=*]
    \item RF10 --- Perfis adicionais: coordenador e magistrado
    \item RF11 --- Verificação de IMEI via API externa
    \item RF12 --- Dashboard analítico por agente e por processo
    \item RF13 --- Criação e associação de despachos judiciais às provas
\end{itemize}

\subsection*{Won't have (nesta versão)}

\begin{itemize}[nosep, leftmargin=*]
    \item RF14 --- Integração com sistemas internos da PSP
                   (requer autorizações não disponíveis)
    \item RF15 --- Dados reais de processos judiciais
                   (questões legais e de confidencialidade)
    \item RF16 --- Módulos de Química, Biologia ou Física Forense
                   (estrutura documentada, não implementada)
    \item RF17 --- Aplicação móvel nativa iOS/Android
                   (interface web \emph{mobile-first} é suficiente para MVP)
\end{itemize}

% =============================================================================
\section{Requisitos não-funcionais}
% =============================================================================

\subsection*{Must have}

\begin{itemize}[nosep, leftmargin=*]
    \item RNF01 --- \textbf{Segurança:} HTTPS obrigatório; autenticação JWT para
          acesso a todos os endpoints protegidos; passwords com hash.
    \item RNF02 --- \textbf{Usabilidade:} Interface \emph{mobile-first} utilizável
          num smartphone no terreno, com uma mão, sem formação técnica prévia.
    \item RNF03 --- \textbf{Integridade:} Hash SHA-256 dos metadados de cada registo
          de prova no momento de criação; log de custódia \emph{append-only}.
\end{itemize}

\subsection*{Should have}

\begin{itemize}[nosep, leftmargin=*]
    \item RNF04 --- \textbf{Performance:} Tempo de resposta inferior a 2 segundos
          para operações principais.
    \item RNF05 --- \textbf{Manutenibilidade:} Código com testes automatizados nos
          módulos críticos (autenticação, custódia).
\end{itemize}

\subsection*{Could have}

\begin{itemize}[nosep, leftmargin=*]
    \item RNF06 --- \textbf{Disponibilidade:} Deploy acessível remotamente para
          demonstração.
\end{itemize}

% =============================================================================
\section{Histórico de alterações}
% =============================================================================

\begin{longtable}{L{2cm} L{3.5cm} L{5.5cm} L{3cm}}
    \toprule
    \textbf{Versão} & \textbf{Data} & \textbf{Alteração} & \textbf{Razão} \\
    \midrule
    \endhead
    1.0 & 22 mar 2026 & Versão inicial & Proposta de projecto \\
    \bottomrule
\end{longtable}

\end{document}
